\documentclass[11pt]{article}

\usepackage[margin=1in]{geometry}
\usepackage{amsmath,amssymb,mathtools}
\usepackage{booktabs}
\usepackage{enumitem}
\usepackage[hidelinks]{hyperref}

\newcommand{\R}{\mathbb{R}}
\newcommand{\Z}{\mathbb{Z}}
\newcommand{\cvar}{\operatorname{CVaR}}
\newcommand{\ind}{\mathbb{I}}

\title{Staged Solution Method for the Type-Indexed Distance-State CVaR Model}
\author{}
\date{}

\begin{document}
\maketitle

\section{Purpose and positioning}

The monolithic stochastic mixed-integer program jointly chooses facility
activations, vehicle bases, scenario-dependent vehicle routes, resource
allocations, and CVaR variables.  This provides the strongest optimization
model, but it places all strategic and operational integer decisions for every
scenario in one branch-and-bound tree.

The staged method is a decomposition-based matheuristic.  It separates the
strategic posture, detailed scenario routing, and final risk-aware allocation:
\begin{equation*}
\begin{gathered}
\boxed{\text{Stage 1: choose }p,b}\\[2mm]
\downarrow\\[2mm]
\boxed{\text{Stage 2: choose }n^\omega\text{ separately by scenario}}\\[2mm]
\downarrow\\[2mm]
\boxed{\text{Stage 3: optimize full-distribution CVaR allocation}}.
\end{gathered}
\end{equation*}

The method produces a feasible solution to the detailed model when all stages
are feasible, but it does not generally provide a globally optimal solution
because decisions fixed in an earlier stage cannot be revised later.

\section{Compact monolithic reference model}

Let
\begin{align*}
s&=(p,b), &&\text{first-stage strategic decisions},\\
v^\omega&=(n^\omega,\bar n^\omega,x^\omega,y^\omega,z^\omega),
&&\text{scenario-$\omega$ recourse decisions}.
\end{align*}
Here $p_i$ activates candidate site $i$, $b_{kj}$ is the number of type-$k$
vehicles based at $j$, and $n^\omega_{kijd}$ is the integer count of type-$k$
vehicles using arc $(i,j)$ after consuming distance state $d$.

Let $\mathcal S$ denote the feasible strategic set containing site-count,
activation-budget, fleet-assignment, and base-linking constraints.  Let
$\mathcal F_\omega(s)$ denote the detailed distance-state recourse set for
scenario $\omega$, including
\begin{itemize}[nosep]
\item distance-expanded vehicle conservation;
\item payload and degraded arc capacity;
\item node-handling capacity;
\item resource balance and unmet demand; and
\item all variable domains.
\end{itemize}

The scenario loss is
\begin{align}
\Lambda^\omega(s,v^\omega)
=
&\sum_{i\in N}\sum_{r\in R}
\delta_{ir}z^\omega_{ir}
\nonumber\\
&+\sum_{k\in K}\sum_{(i,j)\in A_k}\sum_{r\in R}
c_{kij}x^\omega_{kijr}
\nonumber\\
&+\varepsilon
\sum_{k\in K}\sum_{(i,j)\in A_k}
\sum_{d\in\mathcal D_{kij}}n^\omega_{kijd}.
\label{eq:full-loss}
\end{align}

The monolithic reference formulation is
\begin{align}
\min_{s,v,\eta,\xi}\quad
&\eta+
\frac{1}{1-\beta}
\sum_{\omega\in\Omega}\pi^\omega\xi^\omega
\label{eq:monolithic-objective}\\
\text{s.t.}\quad
&s\in\mathcal S,
\label{eq:monolithic-strategic}\\
&v^\omega\in\mathcal F_\omega(s)
&&\forall\omega\in\Omega,
\label{eq:monolithic-recourse}\\
&\xi^\omega\geq
\Lambda^\omega(s,v^\omega)-\eta
&&\forall\omega\in\Omega,
\label{eq:monolithic-cvar}\\
&\xi^\omega\geq0
&&\forall\omega\in\Omega,
\qquad \eta\in\R.
\end{align}

\section{Stage 1A: reduced-scenario detailed strategic solve}

This is the Stage~1 formulation used by the standard staged solver.  Select a
strategic scenario subset $\Omega_S\subseteq\Omega$.  The probabilities used
within the subset are
\begin{equation}
\widehat\pi^\omega
=
\frac{\pi^\omega}
{\sum_{\nu\in\Omega_S}\pi^\nu}
\qquad \forall\omega\in\Omega_S.
\label{eq:renormalized-probability}
\end{equation}

Stage~1A solves
\begin{align}
\min_{s,v,\eta_S,\xi}\quad
&\eta_S+
\frac{1}{1-\beta}
\sum_{\omega\in\Omega_S}\widehat\pi^\omega\xi^\omega
\label{eq:stage1-objective}\\
\text{s.t.}\quad
&s\in\mathcal S,
\label{eq:stage1-strategic}\\
&v^\omega\in\mathcal F_\omega(s)
&&\forall\omega\in\Omega_S,
\label{eq:stage1-recourse}\\
&\xi^\omega\geq
\Lambda^\omega(s,v^\omega)-\eta_S
&&\forall\omega\in\Omega_S,
\label{eq:stage1-cvar}\\
&\xi^\omega\geq0
&&\forall\omega\in\Omega_S,
\qquad \eta_S\in\R.
\end{align}

Let the resulting incumbent strategic posture be
\begin{equation}
s^H=(p^H,b^H).
\label{eq:strategic-incumbent}
\end{equation}
The Stage~1 routes and allocations are used to value $p$ and $b$, but they are
discarded after the solve.  Only $p^H$ and $b^H$ are passed to Stage~2.

If $\Omega_S=\Omega$, Stage~1A contains the complete monolithic formulation.
The subsequent staged procedure nevertheless discards and reconstructs its
operational decisions.

\section{Optional Stage 1B: all-scenario strategic proxy}

The proxy variant removes detailed vehicle-flow variables from Stage~1 and
uses every scenario in $\Omega$.  It retains integer $p$ and $b$ but uses
continuous direct-service variables:
\begin{align*}
t^\omega_{kji}&\geq0,
&&\text{equivalent type-$k$ trips from base $j$ to demand node $i$},\\
u^\omega_{kjir}&\geq0,
&&\text{resource $r$ delivered by those trips},\\
u^{\omega,\mathrm{loc}}_{jr}&\geq0,
&&\text{resource $r$ used locally at activated site $j$},\\
z^\omega_{ir}&\geq0,
&&\text{proxy unmet demand}.
\end{align*}

For each feasible service option $(\omega,k,j,i)$, precompute
\begin{align*}
L^\omega_{kji}&=\text{shortest-path mission distance},\\
C^\omega_{kji}&=\text{minimum of vehicle payload and path bottleneck},\\
\widetilde c^\omega_{kji}&=\text{outbound unit transport cost}.
\end{align*}
An option is omitted if its path is unavailable or
$L^\omega_{kji}>D_k$.

The proxy payload constraints are
\begin{equation}
\sum_{r\in R}w_r u^\omega_{kjir}
\leq
C^\omega_{kji}t^\omega_{kji}
\qquad\forall(\omega,k,j,i)\text{ retained}.
\label{eq:proxy-payload}
\end{equation}

Aggregate fleet-distance capacity is
\begin{equation}
\sum_{i:(\omega,k,j,i)\text{ retained}}
L^\omega_{kji}t^\omega_{kji}
\leq
D_k a^\omega_j b_{kj}
\qquad\forall\omega,k,j.
\label{eq:proxy-distance}
\end{equation}

For one-way proxy missions, each vehicle is used at most once:
\begin{equation}
\sum_{i:(\omega,k,j,i)\text{ retained}}
t^\omega_{kji}
\leq
a^\omega_j b_{kj}
\qquad\forall\omega,k,j.
\label{eq:proxy-one-way-count}
\end{equation}
If return-to-base service is assumed, the return distance is included in
$L^\omega_{kji}$ and constraint~\eqref{eq:proxy-one-way-count} is omitted,
allowing repeated missions within the fleet-distance budget.

Inventory availability is
\begin{align}
u^{\omega,\mathrm{loc}}_{jr}
+\sum_{k\in K}\sum_i u^\omega_{kjir}
&\leq
\rho I_{jr}a^\omega_jp_j
&&\forall\omega,j\in N^P,r.
\label{eq:proxy-inventory}
\end{align}

Demand coverage is
\begin{equation}
\ind_{\{i\in N^P\}}u^{\omega,\mathrm{loc}}_{ir}
+\sum_{k\in K}\sum_{j\in J_k}u^\omega_{kjir}
+z^\omega_{ir}
=d^\omega_{ir}
\qquad\forall\omega,i,r.
\label{eq:proxy-demand}
\end{equation}

The proxy also retains the scenario-degraded node-handling restriction
\begin{equation}
\sum_{k\in K_m}\sum_{j\in J_k}t^\omega_{kji}
\leq
\widehat\Theta^\omega_{i,m}(p)
\qquad\forall\omega,m,i.
\label{eq:proxy-handling}
\end{equation}

Proxy scenario loss is
\begin{align}
\widetilde\Lambda^\omega
=
&\sum_{i\in N}\sum_{r\in R}\delta_{ir}z^\omega_{ir}
\nonumber\\
&+\sum_{k\in K}\sum_{j\in J_k}\sum_{i\in N}\sum_{r\in R}
\widetilde c^\omega_{kji}u^\omega_{kjir}
\nonumber\\
&+\varepsilon\sum_{k\in K}\sum_{j\in J_k}\sum_{i\in N}
t^\omega_{kji}.
\label{eq:proxy-loss}
\end{align}

Stage~1B solves
\begin{align}
\min\quad
&\widetilde\eta+
\frac{1}{1-\beta}
\sum_{\omega\in\Omega}\pi^\omega\widetilde\xi^\omega
\label{eq:proxy-cvar-objective}\\
\text{s.t.}\quad
&s\in\mathcal S,
\quad
\text{constraints~\eqref{eq:proxy-payload}--\eqref{eq:proxy-handling}},
\nonumber\\
&\widetilde\xi^\omega
\geq
\widetilde\Lambda^\omega-\widetilde\eta
&&\forall\omega\in\Omega,
\label{eq:proxy-cvar-excess}\\
&\widetilde\xi^\omega\geq0
&&\forall\omega\in\Omega,
\qquad\widetilde\eta\geq0.
\end{align}

Stage~1B minimizes CVaR of proxy losses, not detailed routed losses.  Its output
is again the strategic incumbent $s^H=(p^H,b^H)$.

\section{Stage 2: independent detailed scenario routing}

Fix the strategic posture obtained from Stage~1:
\begin{equation}
p=p^H,
\qquad
b=b^H.
\label{eq:stage2-fixing}
\end{equation}

For every $\omega\in\Omega$, solve an independent detailed distance-state
routing problem
\begin{align}
Q_\omega(p^H,b^H)
=
\min_{v^\omega}\quad
&\Lambda^\omega(p^H,b^H,v^\omega)
\label{eq:stage2-objective}\\
\text{s.t.}\quad
&v^\omega\in\mathcal F_\omega(p^H,b^H).
\label{eq:stage2-feasibility}
\end{align}

With one scenario, the CVaR expression reduces to that scenario's loss, so
solving the one-scenario CVaR model is equivalent to
equations~\eqref{eq:stage2-objective}--\eqref{eq:stage2-feasibility}.

Let the selected integer route be
\begin{equation}
n^{H,\omega}
\qquad\forall\omega\in\Omega.
\label{eq:stage2-routes}
\end{equation}
The resource-allocation variables used to discover this route are discarded;
they are reoptimized jointly in Stage~3.  The scenario problems are mutually
independent once $p^H$ and $b^H$ are fixed and may therefore be solved in
parallel.

\section{Stage 3: full-distribution CVaR allocation LP}

Fix all strategic and route decisions:
\begin{equation}
p=p^H,
\qquad
b=b^H,
\qquad
n^\omega=n^{H,\omega}
\quad\forall\omega\in\Omega.
\label{eq:stage3-fixing}
\end{equation}

The remaining decision variables are continuous resource allocations,
terminations, unmet demand, retained resources, and CVaR variables.  Define
$\mathcal X_\omega(p^H,b^H,n^{H,\omega})$ as the continuous portion of the
detailed recourse set.  Stage~3 solves
\begin{align}
z^H=\min_{x,y,z,\bar n,\eta,\xi}\quad
&\eta+
\frac{1}{1-\beta}
\sum_{\omega\in\Omega}\pi^\omega\xi^\omega
\label{eq:stage3-objective}\\
\text{s.t.}\quad
&(x^\omega,y^\omega,z^\omega,\bar n^\omega)
\in
\mathcal X_\omega(p^H,b^H,n^{H,\omega})
&&\forall\omega\in\Omega,
\label{eq:stage3-allocation}\\
&\xi^\omega\geq
\Lambda^\omega(p^H,b^H,n^{H,\omega},x^\omega,y^\omega,z^\omega)-\eta
&&\forall\omega\in\Omega,
\label{eq:stage3-cvar}\\
&\xi^\omega\geq0
&&\forall\omega\in\Omega,
\qquad\eta\in\R.
\end{align}

Because $p$, $b$, and $n$ are fixed, Stage~3 is a linear program.  Its value
$z^H$ is the best full-distribution CVaR attainable with the staged posture and
routes.

\section{Complete staged algorithm}

\begin{enumerate}[label=\textbf{Step \arabic*:},leftmargin=18mm]
\item Generate the common scenario set $\Omega$, probabilities $\pi^\omega$,
      demands, severities, and degraded capacities.
\item Solve either Stage~1A or Stage~1B and retain $p^H,b^H$.
\item For each $\omega\in\Omega$, fix $p^H,b^H$, solve the detailed routing
      problem, and retain $n^{H,\omega}$.
\item Build the all-scenario model, fix $p^H,b^H,n^{H,\omega}$, and solve the
      continuous CVaR allocation problem.
\item Report $z^H$, $\eta^H$, scenario losses, tail scenarios, selected sites,
      vehicle base counts, routing gaps, and total runtime.
\end{enumerate}

\section{Feasibility, bounds, and interpretation}

If every Stage~2 route is feasible and Stage~3 is feasible, the assembled
solution is feasible for the detailed monolithic formulation.  Since the
problem is a minimization problem,
\begin{equation}
z^*\leq z^H,
\label{eq:upper-bound}
\end{equation}
where $z^*$ is the unknown monolithic optimum.  Thus $z^H$ is a valid upper
bound for the same detailed formulation and scenario distribution.

The staged method does not generate a global lower bound by itself.  If $LB$
is obtained from a valid LP relaxation or a monolithic branch-and-bound run,
the staged solution may be assessed using
\begin{equation}
\operatorname{Gap}_H
=
\frac{z^H-LB}{\lvert z^H\rvert}\times100\%.
\label{eq:heuristic-gap}
\end{equation}

Stage~3 is conditionally optimal:
\begin{equation}
z^H
=
\min\left\{
\cvar_\beta(\Lambda^\omega):
p=p^H,\ b=b^H,\ n^\omega=n^{H,\omega}
\right\}.
\end{equation}
It is not generally globally optimal over all possible $p$, $b$, and $n$.

\section{Comparison with the monolithic method}

\begin{center}
\begin{tabular}{p{0.28\textwidth}p{0.31\textwidth}p{0.31\textwidth}}
\toprule
Feature & Monolithic method & Staged method \\
\midrule
Strategic and operational decisions
& Optimized jointly
& Fixed sequentially \\
Scenario routing
& One extensive-form MIP
& One independent MIP per scenario \\
Final CVaR
& Jointly determines all decisions
& Reoptimizes allocation conditional on fixed decisions \\
Global lower bound
& Produced by branch-and-bound
& Requires a separate relaxation or monolithic run \\
Global optimality
& Possible if solved to proof
& Not guaranteed \\
Parallel opportunity
& Primarily solver-internal
& Stage~2 scenarios are independent \\
Operational output
& Feasible incumbent if found
& Feasible assembled solution if all stages succeed \\
\bottomrule
\end{tabular}
\end{center}

\section{Recommended reporting language}

The method should be described as a \emph{staged decomposition-based
matheuristic} or \emph{sequential fix-and-optimize method}.  It should not be
described as Benders decomposition because no feasibility or optimality cuts
are returned to a strategic master problem.  It should also not be described
as exact unless the full monolithic formulation is solved to a valid global
optimality certificate.

For a fair comparison between staged variants, compare the final Stage~3
objective values under the same scenarios, probabilities, $\beta$, distance
buckets, and detailed routing formulation.  Stage~1A and Stage~1B objective
values are not directly comparable because they use different recourse models
and may use different scenario sets.

\end{document}
